% STATUS: DRAFT
\chapter{\cstar-algebras}\label{ch:cstar}

\begin{physmotiv}
What do we know about the observables of a system before we choose a state or a Hilbert space? That they can be added, multiplied by complex numbers, multiplied together (non-commutatively), conjugated ($A\mapsto A\ad$), and that they have a norm (a measure of the largest possible measured value). A \cstar-algebra is exactly this: nothing more, nothing less. For a system with finitely many degrees of freedom the algebra is $\BH$; for a spin chain or a field one has to construct it directly from the local commutation relations, since no preferred $\HH$ exists (\cref{ch:why_algebras}).
\end{physmotiv}

\section{Definition and the basic theorem}

\begin{definition}
A \emph{\cstar-algebra} is a complex algebra $\A$ with an involution $a\mapsto a\ad$ and a norm $\norm\cdot$ such that $\A$ is complete, $\norm{ab}\le\norm a\norm b$, and
\begin{equation}
\norm{a\ad a}=\norm a^2\qquad(\text{the \cstar-identity}).
\end{equation}
We take $\A$ unital (with unit $\id$) unless stated otherwise.
\end{definition}

The \cstar-identity ties the norm to the algebra. It implies that the norm is determined by the algebraic structure: $\norm a^2$ is the spectral radius of $a\ad a$. Hence an algebra has \emph{at most one} \cstar-norm, and every $^*$-homomorphism between \cstar-algebras is automatically norm-decreasing.

\begin{theorem}[Gelfand--Naimark]
Every \cstar-algebra is isometrically $^*$-isomorphic to a norm-closed $^*$-subalgebra of $\Bnd{\HH}$ for some Hilbert space $\HH$.
\end{theorem}
Thus abstract \cstar-algebras and norm-closed algebras of operators are the same thing, but the Hilbert space is not unique. The construction of such a $\HH$ from a state is the GNS construction (\cref{ch:gns}).

\subsection*{Spectrum, positivity, functional calculus}
An element $a$ is self-adjoint if $a=a\ad$. The \emph{spectrum} $\spec(a)=\{\lambda\in\C:a-\lambda\id\text{ not invertible}\}$ is a non-empty compact set, real for self-adjoint $a$, and \emph{independent of which algebra} (or representation) it is computed in: this is crucial since it means $\spec(a)$ is intrinsic to $\A$. For self-adjoint $a$ one has the \emph{continuous functional calculus}: $f(a)\in\A$ for every continuous function $f$ on $\spec(a)$. An element is \emph{positive} ($a\ge0$) if it is self-adjoint with $\spec(a)\subset[0,\infty)$; equivalently $a=b\ad b$ for some $b$. Positivity gives the order structure on observables used to define states (\cref{ch:states}).

\begin{whatwrong}
Only \emph{continuous} functions of $a$ lie in the \cstar-algebra. Spectral projections $\chi_\Omega(a)$ (functions with jumps), in general do \emph{not}: take $\A=C[0,1]$ and $a(x)=x$; then $\chi_{[0,1/2]}(a)$ is discontinuous and not in $C[0,1]$. They do lie in the von Neumann algebra obtained by weak closure (\cref{ch:doublecomm}). This is a reason one passes from \cstar\ to vN algebras: yes/no questions need to be elements.
\end{whatwrong}

\section{Examples}

\begin{enumerate}
\item \textbf{Matrices} $M_n(\C)=\Bnd{\C^n}$. The universal finite-dimensional example; every finite-dimensional \cstar-algebra is a direct sum $\bigoplus_kM_{n_k}(\C)$.
\item \textbf{Compact operators} $\mathcal K(\HH)$ (non-unital) and $\BH$.
\item \textbf{Commutative algebras} $C(X)$, continuous complex functions on a compact Hausdorff space $X$ with the sup-norm. By the Gelfand--Naimark \emph{duality} (\cref{ch:gelfand}) every commutative unital \cstar-algebra has this form with $X$ the space of characters. Commutative \cstar-algebras are classical phase spaces.
\item \textbf{Weyl (CCR) algebra} $\mathcal W(\mathcal S,\sigma)$ of a symplectic space: the \cstar-algebra generated by unitaries $W(f)$ with $W(f)W(g)=\ee^{\frac\ii2\sigma(f,g)}W(f+g)$ (\cref{ch:weyl}). It has a unique \cstar-norm (Slawny), so it is a well-defined object independent of representation. For a free scalar field of mass $m$ in Minkowski space, $\sigma$ is the Pauli--Jordan commutator and the algebra of a region $R$ is the subalgebra $\mathcal W(\mathcal S_R)$ generated by $W(f)$ with $f$ supported in $R$.
\item \textbf{CAR algebra} of fermions: generated by $a(f)$, $f\in\mathcal K$ a Hilbert space, with $\{a(f),a(g)\ad\}=\inner fg\id$, $\{a(f),a(g)\}=0$. Since $a(f)$ is automatically bounded ($\norm{a(f)}=\norm f$), no Weyl trick is needed.
\item \textbf{The infinite spin chain (UHF algebra).} Let $\A_N=\bigotimes_{j=1}^NM_2(\C)=M_{2^N}(\C)$, embedded as $\A_N\hookrightarrow\A_{N+1}$ via $a\mapsto a\otimes\id$. The norm-closure of $\bigcup_N\A_N$ is the \cstar-algebra $\A=M_{2^\infty}$ of the infinite chain. It is \emph{simple} (no non-trivial closed two-sided ideals) and unique: one algebra, an infinite family of different representations (\cref{ch:inequiv}), as in \cref{ch:why_algebras}. It carries a unique \emph{tracial} state $\tau(ab)=\tau(ba)$, the infinite-temperature state, which we will meet again as the origin of the hyperfinite type II$_1$ factor (\cref{ch:classification}).
\end{enumerate}

\begin{keyidea}
A \cstar-algebra is the representation-independent algebra of observables. Matrices, classical phase spaces (commutative case), CCR/CAR algebras, and spin-chain algebras are all \cstar-algebras. Continuity of the norm means only \emph{continuous} functions of observables are in it.
\end{keyidea}

\section{Locality and the quasi-local picture}

Algebras of \emph{local observables}, those supported in a finite region $\Lambda$, are $\A(\Lambda)$ for each finite $\Lambda$ and one has isotony $\Lambda_1\subset\Lambda_2\Rightarrow\A(\Lambda_1)\subset\A(\Lambda_2)$, and commutation $[\A(\Lambda_1),\A(\Lambda_2)]=0$ if $\Lambda_1\cap\Lambda_2=\emptyset$ (for bosons/spins; graded for fermions). The \emph{quasi-local algebra} is the norm closure $\A=\overline{\bigcup_\Lambda\A(\Lambda)}$. This is the infinite spin chain example and the framework of \cref{ch:quasilocal}. In relativistic QFT we replace $\Lambda$ by spacetime regions (\cref{ch:haagkastler}).

\section*{Exercises}
\begin{exercise}
Show from the \cstar-identity that $\norm{a\ad}=\norm a$ and that for self-adjoint $a$, $\norm{a^{2^k}}=\norm a^{2^k}$, hence $\norm a=\lim\norm{a^n}^{1/n}$ (spectral radius). Conclude the norm is determined by the algebra.
\end{exercise}
\begin{exercise}
Show that $\A_N=M_{2^N}$ and the map $a\mapsto a\otimes\id$ is an isometry, so that the inductive limit norm is well defined. Show that $\sigma^z_j$ and $\sigma^x_k$ commute for $j\ne k$. Is $M_N=\frac1N\sum_{j\le N}\sigma_j^z$ in $\A$? Does the sequence converge in norm? (Compare \cref{ch:topologies}.)
\end{exercise}
\begin{exercise}
For $\A=C[0,1]$, find the spectrum of $a(x)=x$, compute $\norm a$, and show $\chi_{[0,1/2]}(a)\notin\A$. Show that every closed two-sided ideal of $C(X)$ is of the form $\{f:f|_Y=0\}$ for a closed $Y\subset X$.
\end{exercise}
\begin{exercise}
Show that for fermionic modes $a(f)$ one has $\norm{a(f)}=\norm f$ using $a(f)a(f)\ad+a(f)\ad a(f)=\norm f^2\id$. Explain why one can treat fermions without exponentiating.
\end{exercise}
